\section{Philokalia and the Recovery of Tradition}

\begin{quotex}
The Kingdom of Heaven suffers violence and the violent take it by force. (Matthew 11:12)
\end{quotex}


There is a new edition of Volume I of the Philokalia\footnote{\url{http://www.ibmgs.org/patristic.html}} translated by Constantine Cavarnos and published by the Institute for Byzantine and Modern Greek Studies\footnote{\url{http://www.ibmgs.org}}. \emph{Philokalia} means ``Love of the Beautiful'' (not ``beautiful things'') and serves as a guidebook for early Christian spiritual practices. Besides the Bible, the works included in the Philokalia are heavily influenced by Platonic notions; it would even be reasonable to speculate that, as such, it includes spiritual practices that were part of the training at Plato's Academy. This work is recommended as essential to the recovery of the ``metaphysical view prior to the scholastic age\footnote{\url{https://gornahoor.net/?p=488}}''.

The differences between the Philokalia and contemporary manifestations of Christianity are striking.

\begin{itemize}


\item \textbf{Salvation}. Salvation is the process of \emph{theosis}, or divinization, that is, becoming like God. This differs from both from

  \begin{itemize}

\item the Roman notion that salvation comes from following a set of moral norms and seeking reconciliation after sinning, and
  
\item the Protestant notion of imputed righteousness devoid of any real, actual, and effective change.
  \end{itemize}


\item \textbf{Effort}. \emph{Theosis} requires conscious efforts, both physical and spiritual. Some of the latter include concentration, meditation, inner attention, and mental prayer. All these involve control of thought so the mind doesn't wander aimlessly\footnote{\url{https://gornahoor.net/?p=451}}.

\item \textbf{Rational}. The necessary first step is to be ``rational'', that is, in harmony with the \emph{Logos}. Hence, irrationality, ungrounded exuberance, and blind faith are obstacles on the path.

\item \textbf{Transformation}. The end result is the transformation of one's level of being\footnote{\url{https://gornahoor.net/?p=478}}. Every man is in the image of God, but few are in the likeness of God. According to St Anthony, the former are barely human and should be avoided.
\end{itemize}


Although these writings are geared to the contemplative life (Guenon), we agree with Evola that the way of action is a valid path for certain personalities. Hence, some adaptations need to be made in the latter case. In more contemporary terms, we may call one path the Way of the Monk, and the other, the Way of the Man in the World.\footnote{The first essay in the Philokalia is attributed to St Anthony and concerns the attributes and qualities of a rational man. In the first English edition translated by Palmer, Sherrard, and Ware, this essay is removed from its prime position and relegated to a mere appendix on the grounds it is insufficiently ``Christian''. I'll not bother with their reasons, since all this demonstrates is the incomprehensibility of early Christianity to our contemporaries, who even claim to be Christian leaders. Let this serve as a warning to anyone seeking authentic Christian doctrine and spirituality. What the first translators have done is to split the Western tradition in two, and detach the Christian religion from the Primordial Tradition.}



\flrightit{Posted on 2010-05-25 by Cologero}

\begin{center}* * *\end{center}

\begin{footnotesize}
\begin{sffamily}
\texttt{Matt on 2010-05-26 at 14:27 said:}

Is the St. Antony mentioned Anthony the Hermit, Anthony of Kiev, or Antony of Siya?

\hfill

\texttt{Cologero on 2010-05-26 at 23:21 said:}

St Antony the Great aka St Antony of Egypt

\hfill

\texttt{Matt on 2010-05-29 at 15:43 said:}

Thank you.

\hfill

\texttt{James O'Meara on 2010-12-11 at 09:01 said:}

FWIW, I have several of those ``Skylight Illuminations'' texts that print notes and commentary on the facing pages. For the selection of texts from the Philokalia, there are two paragraphs from Antony, sections 66 and 79 of his text, under the chapter headed ``The Passions''. The accompanying note says the work is ``actually of non-christian origins, reflecting platonic and stoic teachings'' lower case theirs , which just about makes your point.

\hfill

\texttt{Matt on 2010-12-11 at 11:39 said:}

James,

My school library has those translations as well. I think whats even funnier is that one of the translators, Philip Sherrard, is considered a figure among perennialism... ..yet he was perfectly fine with devaluing an important writing of early Christianity due to attributes of Platonic and Stoic thinking in it. Very perennial of him.

\hfill

\texttt{James O'Meara on 2010-12-12 at 10:15 said:}

Matt,

PS had an article back in the 80s in some `scholarly' philosophical journal, called `Tradition and the Traditions' which attempted a `logical' argument precisely to `devalue' the Platonic element in Christianity, and more broadly, Traditionalism `what you mean by `higher' or `more logical' is relative to your tradition' etc. Very odd; trying to gain favor with The Establishment?

Now, in the new edition of Schuon's Logic and Transcendence, there are some `unpublished letters' added, and lo, there is a passage that must refer to PS: ``This may not be to the liking of S. his or the editor's abbreviation , who persists in his perfectly absurd anti-Platonism in his new article; it is a veritable perversion of intelligence.'' Based on what he says there and in the main text, I must say I side with Schuon!

\end{sffamily}
\end{footnotesize}

